\subsection{Further robustness and implementation requirements}

Besides shifting the interval edges, the economic-parameter perturbations and the historical volatility anchors also moved the critical weight, as Table~\ref{tab:critsens} reports. In every positive-radius case, the critical weight stayed above one half, and the member-threshold ordering was unchanged.

\begin{table}[!t]
\caption{Sensitivity of the critical weight under the volatility-only ambiguity declaration.}
\label{tab:critsens}
\centering
\scriptsize
\setlength{\tabcolsep}{3pt}
\begin{tabular}{@{}lcccc@{}}
\toprule
\multicolumn{5}{@{}l}{\emph{Panel A. Economic-parameter sensitivity at a 10\% relative radius}}\\
Perturbation & \multicolumn{4}{r}{$\tilde{\alpha}$ at 10\%}\\
\midrule
Baseline calibration & \multicolumn{4}{r}{0.517}\\
$r=0.06$ ($-25\%$) & \multicolumn{4}{r}{0.520}\\
$r=0.12$ ($+50\%$) & \multicolumn{4}{r}{0.514}\\
$\mu_c=0.06$ ($+50\%$) & \multicolumn{4}{r}{0.521}\\
$\sbar=0.199$ ($-20\%$) & \multicolumn{4}{r}{0.517}\\
$\sbar=0.299$ ($+20\%$) & \multicolumn{4}{r}{0.518}\\
\midrule
\multicolumn{5}{@{}l}{\emph{Panel B. Historical-anchor sensitivity across three relative radii}}\\
Historical anchor regime & Anchor estimate & 10\% & 22\% & 33\%\\
\midrule
\makecell[l]{Shanghai Emission Allowance\\(SHEA) pilot, 2016 to 2021Q1} & $\sbar=0.953$ & 0.524 & 0.553 & 0.580\\
Early CEA, 2021Q3 to 2023Q4 & $\sbar=0.299$ & 0.518 & 0.540 & 0.560\\
\makecell[l]{Current CEA, 2024 to 2025\\(baseline)} & $\sbar=0.249$ & 0.517 & 0.538 & 0.557\\
\makecell[l]{Pooled national,\\2021Q3 to 2025Q4} & $\sbar=0.278$ & 0.518 & 0.539 & 0.558\\
\bottomrule
\end{tabular}
\par\smallskip
\parbox{\columnwidth}{\scriptsize\emph{Notes.} The $\sbar=0.299$ perturbation of Panel A coincides numerically with the early CEA anchor regime of Panel B.}
\end{table}

\paragraph{Implementation safeguards} An elicitation instrument for the recorded weights must meet three requirements. It must return a number in $[0,1]$ on the adverse endpoint of the declared ambiguity set. A calibration to a different object needs a stated mapping. It must be reproducible at the same gate, because the recorded weight enters the gate audit record. Its precision must also be assessed against the review band around the critical weight. A prospective implementation should therefore prespecify standardized administration with attention and comprehension checks. It should also prespecify a test-retest assessment of rank-order reliability and a check of scale comparability across gates. Candidate instruments include matching-probability tasks in the style of \citet{Dimmock2016}, lottery equivalents, and structured vignettes. A discrete-choice format such as best-worst scaling could also serve, with its output mapped onto $[0,1]$ by a rule fixed in advance. The instrument version, acceptance rule, and recorded weights should be frozen before later gate outcomes are observed. An extreme weight sets an interval edge, so it should prompt a documented review before any exclusion. Two further choices belong to institutional policy, and the first is reporting granularity. The edges, the width, and the three-state summary are symmetric functions of the member thresholds, so a sorted list without member identities preserves them. A sorted list, however, removes the link between each threshold and its member. An identified table keeps this link, although it also exposes the most divergent member in the minutes, as the documented review of extreme weights does. The second choice is the elicitation schedule, which matters because the extreme recorded weights alone determine both edges. A member aware of this mapping could move the reported interval through the choice of weight, so the method assumes weights recorded in advance. Collecting the weights before project materials circulate establishes a pre-disclosure reference, provided the elicitation names the common information behind them. A project-specific adjustment of the weights would require a second elicitation under a documented information set.

The acceptance rule for the instrument can rest on a precision target derived from the model. Under the classical error model, measurement noise with standard deviation $\sigma_{\alpha}$ enters each recorded weight. A member whose recorded weight lies a distance $d$ from the critical weight then receives the opposite directional label with probability $\Phi(-d/\sigma_{\alpha})$. Holding this probability at or below $p$ requires $\sigma_{\alpha}\le d/\Phi^{-1}(1-p)$. With $d$ set to the illustrative review band of 0.05, the condition gives $\sigma_{\alpha}\le0.039$ for a 10\% flip probability and $\sigma_{\alpha}\le0.030$ for 5\%. These values serve as design targets for the ordinal use of the weights. Members closer than $d$ to the critical weight keep de-emphasized labels, since no instrument precision secures a stable label arbitrarily close to it. Beyond noise within a gate, recorded weights can also drift between gates. When the recorded weights of simulated panels drifted by 0.05 per gate, the three-state summary changed in 28.9\% of the panels. The re-elicitation cadence is therefore a design parameter.

These elicitation choices join the other configuration decisions a prospective gate must record, which Table~\ref{tab:config} organizes by the three layers of Figure~\ref{fig:architecture}. The text-pipeline items and the constant $c_w$ serve only a radius derived from documents, so a committee declaring its radii directly drops them. Its threshold report stays reproducible from the recorded radius, although the expert reasoning behind this radius cannot be regenerated from a text pipeline. In the interface layer, the registry snapshot authorizes affine-only release, so an institution always issuing the dual report needs no registry. It then loses both the option to omit the benchmark and the record of the configuration authorizing a past release. It can still retain and replay the inputs and outputs of each gate.

The irreducible configuration is shorter than Table~\ref{tab:config} suggests. At minimum, a prospective configuration fixes four items: the economic primitives, the volatility anchor, the declared coordinate, and the weight-recording protocol. Each gate then supplies its declared radius, the recorded member weights, and the prevailing price. Benchmark recomputation also requires the registered reference project value. The declared coordinate cannot be dropped, because the directional labels are unreadable without it. Where no member weights are recorded, a hypothetical profile still yields a well-defined interval, which serves as a sensitivity band over a stipulated grid.

\begin{table*}[!t]
\caption{Configuration record required before a prospective gate, organized by the three layers of Figure~\ref{fig:architecture}.}
\label{tab:config}
\centering
\footnotesize
\renewcommand{\arraystretch}{1.1}
\begin{tabular}{@{}L{0.16\textwidth}L{0.3\textwidth}L{0.2\textwidth}L{0.26\textwidth}@{}}
\toprule
Configuration item & Registered before the gate & Opens a new dated cell & Left to the deploying institution\\
\midrule
\multicolumn{4}{@{}l}{\emph{Input layer}}\\
Volatility anchor $\sbar$ & Estimation window, market, and estimator & New window, market boundary, or estimator & An option-implied or otherwise forward-looking anchor\\
Signal-to-radius constant $c_w$ & Target radius $h_{\mathrm{ref}}$, reference signal $w_{\mathrm{ref}}$, and $c_w=h_{\mathrm{ref}}/w_{\mathrm{ref}}$ & Any change to either input, or to the encoder & An empirical calibration, unidentified by the available price record\\
Text pipeline & Encoder, layer, pooling, normalization, truncation, corpus screen & Any change to any of these items & Evidence of semantic dispersion predicting realized volatility\\
Recorded weights $\{\alpha_i\}$ & Instrument version, mapping interval, review band, and acceptance rule. The rule must meet the precision target $\sigma_{\alpha}\le0.039$ at $d=0.05$ and $p=0.10$ & New instrument or mapping interval & Universal acceptance cutoffs and a committee pilot\\
\multicolumn{4}{@{}l}{\emph{Analytical layer}}\\
Economic primitives & $r$, $\mu_c$, $I$, $Q$, and the admissibility margin $\epsilon$ & Any change & Audited full-cost components for a viability appraisal\\
Reference value $V_0/I$ & The value at which the benchmark is evaluated & Any change & Release authorization at $V_0/I\neq1$\\
Declared coordinate & Standard deviation or variance, named explicitly & Any change & None\\
\multicolumn{4}{@{}l}{\emph{Interface layer}}\\
Registry snapshot & Version $\nu$, canonical serialized keys, envelope version, solver version & Any change to any identifier & A continuous-radius or global error certificate\\
Release and reporting & Review cadence, dual-report format, directional-label conventions & Schedule or format change & A voting rule, a consensus procedure, or a full project appraisal\\
\bottomrule
\end{tabular}
\par\smallskip
\parbox{0.97\textwidth}{\footnotesize\emph{Notes.} Every item is frozen before later gate outcomes are observed. Release matching uses the frozen serialized key from the registered pipeline. The fourth column lists the items a deploying institution supplies.}
\end{table*}

\subsection{Preconditions on the reported basis}

Beyond the recorded configuration, five preconditions govern whether the gate can accept the basis of the reported interval. Each precondition names an accounting choice, a model assumption, or an input-validity question, and each is documented before the report is interpreted.

The first precondition is the cost basis, because the reported threshold is marginal. A committee comparing it with a figure including operating costs first applies the common adjustment. Across the 18 national ETS gates with archived prices, no net cost changed the reported classification, because these prices already lay below the interval. A net subsidy changed the classification only at tens of CNY/t, a scale specific to the reference configuration. One gate turned mixed at $\Delta=-10$\,CNY/t. At $\Delta=-50$\,CNY/t, ten of the 18 gates read unanimous invest, and eight remained unanimous wait. At $\Delta=-80$\,CNY/t, all 18 read unanimous invest. The reported thresholds also carry basis risk against the realized carbon cost of a plant. Two common cost components map onto the adjustments of the reported basis. A per-tonne transport and storage charge is the additive adjustment $\Delta$. An energy penalty scales the captured volume by a common factor at fixed utilization, so it is a multiplicative adjustment.

The second precondition is the horizon. Replacing the perpetuity with a finite-horizon annuity multiplies every member threshold, the interval width, and every price-unit envelope by one common factor. The member ordering, the directional labels, and the relative spread survive, while the classification of a given price may change. A gate reporting on an annuity basis therefore rescales the envelope with the thresholds. A finite horizon can also enter the stopping problem. Between a five-year horizon and the perpetuity, it moved the critical weight by less than 0.03 at every examined radius.

The third precondition is the ambiguity channel, since the model places ambiguity on volatility alone and holds drift fixed. At a 10\% volatility radius, adding drift ambiguity of the same relative size multiplied the interval width by 2.41. It also raised the maximum relative trigger error from 0.26\% to 1.48\%. This joint value at a 10\% radius already exceeded the 1.26\% maximum of the volatility-only declaration at radii up to 22\%. A gate under a joint declaration therefore re-evaluates accuracy before release.

The fourth precondition is the price process. On the 1,073 archived log prices, an augmented Dickey-Fuller regression did not reject a unit root, with or without a trend. The Lo-MacKinlay variance-ratio test did not reject a random walk at two, five, ten, or twenty days. These tests did not reject the persistence assumed by the geometric Brownian motion for the price level. Gaussian increments, however, were rejected for the returns of the same series. The implied mean-reversion speed was 0.818 per year, but its confidence interval contained zero, so the series did not identify a speed. The point estimate still lay above the largest reversion speed in Table~\ref{tab:departures}. On this grid, a positive speed moved the anchor threshold from 116.88\,CNY/t to between 97.16 and 99.81\,CNY/t. The anchor threshold of 116.88\,CNY/t was recomputed on the same Crank-Nicolson grid, so the comparison isolates mean reversion from the solver's level bias. The level shift alone exceeded the 7.73\,CNY/t reference margin, so a gate expecting mean reversion should recalibrate. Moving half the declared variance into jumps, at the same total, had a smaller effect. The anchor threshold moved from 116.88 to 115.41\,CNY/t, the interval width changed by at most 1.04\%, and the critical weight rose.

The fifth precondition is the radius input, because the mapping of the policy-text series to realized volatility was unstable in the available record. A registered constant therefore sets the radius scale, and the reference result is sensitive to this constant. An alternative radius based on the sampling uncertainty of the volatility anchor was 6.37\%, close to the 6.057\% peak of the reference path. It left the archived price outside the interval. A radius based on the dispersion of quarterly realized volatility was about seven times larger and made the gate mixed. Every evaluated input alternative changing the classification lay outside the 10\% cap.

Outside the CCUS setting, a transfer to a new application takes six steps, and the first three rebuild the theory. The first defines a common ordered uncertainty set and the meaning of the recorded weight on its adverse endpoint. The second derives the stopping object and verifies the existence of a unique scalar boundary. The third identifies the adverse endpoint as the endpoint with the lower trigger and then rederives trigger monotonicity, curvature, and the critical weight. Three further steps rebuild the reporting apparatus. The fourth maps the member triggers to a common operational scale. The fifth reconstructs the finite benchmark comparison and the versioned release registry. The sixth validates the input, the elicitation, and the scale before deployment. The workflow carries over across applications, whereas the formulas and signs obtained here need rederivation or checking.

The CCUS reconstruction points to three patterns. First, the reference configuration kept every gate at unanimous wait, whereas changing one threshold-side input could produce binding configurations. Second, the common appraisal inputs moved the interval edges more than member disagreement did. Third, the five preconditions on the reported basis shaped the reading of the thresholds. These patterns come from one retrospective reconstruction and may differ in other projects. They nonetheless indicate where a committee might focus before reading the member-threshold vector.
